% evaluation_methodology.tex — GATeR vs TARGET: Evaluation Methodology & Scripts Guide
% Compile: pdflatex evaluation_methodology.tex
\documentclass[11pt,a4paper]{article}
\usepackage[utf8]{inputenc}
\usepackage[T1]{fontenc}
\usepackage[margin=1in]{geometry}
\usepackage{booktabs}
\usepackage{graphicx}
\usepackage{hyperref}
\usepackage{listings}
\usepackage{xcolor}
\usepackage{amsmath}
\usepackage{caption}
\usepackage{enumitem}
\usepackage{multirow}
\usepackage{array}

\definecolor{codebg}{rgb}{0.95,0.95,0.97}
\definecolor{codeframe}{rgb}{0.6,0.6,0.7}
\lstset{
    backgroundcolor=\color{codebg},
    frame=single,
    rulecolor=\color{codeframe},
    basicstyle=\ttfamily\small,
    breaklines=true,
    columns=fullflexible,
    keepspaces=true,
    showstringspaces=false,
}

\title{\textbf{GATeR vs TARGET: Evaluation Methodology}\\[0.5em]
\large Fair Comparison of Graph-Aware RAG Repair vs.\ Fine-Tuned Seq2Seq Repair\\
on the TaRBench Dataset}
\author{GATeR Project --- FYP Defense Reference Document}
\date{\today}

\begin{document}
\maketitle

% ══════════════════════════════════════════════════════════════════════════════
\section{Overview}
% ══════════════════════════════════════════════════════════════════════════════

This document describes the evaluation methodology used to compare
\textbf{GATeR} (Graph-Aware Test Repair) and \textbf{TARGET}
(TaRGet: A Supervised Fine-Tuned Model for Test Repair)
on the same subset of the TaRBench test split (7,103 Java test-repair cases).

\begin{itemize}
    \item \textbf{GATeR}: A zero-shot RAG (Retrieval-Augmented Generation)
          pipeline that builds a \emph{Knowledge Graph} over the project,
          scores context with the novel \emph{KGCompass} formula, compresses
          and aggregates the context, and delegates repair generation to
          DeepSeek Coder 6.7B via Ollama.
    \item \textbf{TARGET}: A supervised model that fine-tunes CodeT5+ (770M)
          on 36,639 TaRBench training pairs and uses beam search ($k{=}40$)
          at inference.
\end{itemize}

\textbf{Fundamental architectural difference.}
TARGET is a \emph{saved model checkpoint}---a file you load and call
\texttt{model.generate()}.
GATeR is a \emph{live pipeline} that queries a Knowledge Graph, retrieves
context from a vector store, compresses and aggregates that context, and
then calls an LLM API to produce the repair.
During evaluation, TARGET's predictions are \textbf{pre-computed} (just loaded
from JSON), whereas GATeR must run its \textbf{full 9-step pipeline} for every
test case.

% ══════════════════════════════════════════════════════════════════════════════
\section{How Each System Produces a Repair}
\label{sec:systems}
% ══════════════════════════════════════════════════════════════════════════════

\subsection{TARGET (No LLM at Evaluation Time)}

TARGET already ran inference on the full TaRBench test set.
Its results are stored in two JSON files:

\begin{enumerate}
    \item \texttt{test\_predictions.json} --- 7,103 entries, each with
          a \texttt{target} (ground truth) and \texttt{preds} (list of
          40 beam candidates).
    \item \texttt{test\_verdicts.json} --- 284,120 entries recording
          whether each beam candidate actually compiled and passed the test.
\end{enumerate}

\textbf{To evaluate TARGET we simply load these files.}
No model inference, no GPU, no LLM server.

\subsection{GATeR (Full Pipeline at Evaluation Time)}

For each test case the evaluation script calls
\texttt{GATREngine.repair\_test()}, which executes the complete
GATeR pipeline:

\begin{enumerate}
    \item \textbf{Step 1 --- Raw Context Ingestion.}
          Query the project's Knowledge Graph (Kuzu) for related entities
          and the Vector Store (LanceDB) for semantically similar code.
    \item \textbf{Steps 2.1--2.6 --- Context Compression.}
          Hybrid scoring (KGCompass $\times$ 0.7 + semantic $\times$ 0.3),
          threshold filtering, snippet compression, relationship filtering,
          pattern extraction, semantic examples.
    \item \textbf{Steps 3.1--3.4 --- RAG Aggregation.}
          Entity aggregation, API delta extraction, canonical usage
          identification, repair strategy formulation.
    \item \textbf{Step 4 --- LLM Repair Generation.}
          The compressed, aggregated context is formatted into a structured
          prompt and sent to DeepSeek Coder 6.7B via Ollama.
          The LLM produces the repaired code.
\end{enumerate}

This means GATeR's evaluation \textbf{requires Ollama running} with the
model loaded, and each case takes 10--60 seconds (KG query + context
compression + LLM generation).

\begin{table}[h]
\centering
\caption{What each system needs during evaluation}
\renewcommand{\arraystretch}{1.15}
\begin{tabular}{@{}lcc@{}}
\toprule
\textbf{Requirement} & \textbf{TARGET} & \textbf{GATeR} \\
\midrule
Running LLM server & No & Yes (Ollama) \\
GPU recommended & No & Yes \\
Knowledge Graph (Kuzu) & No & Yes (per project) \\
Vector Store (LanceDB) & No & Yes (per project) \\
Time per case & $\sim$0\,s (load JSON) & $\sim$10--60\,s \\
Total for 500 cases & $<$1 min & $\sim$2--8 hours \\
\bottomrule
\end{tabular}
\end{table}

% ══════════════════════════════════════════════════════════════════════════════
\section{Evaluation Scripts}
\label{sec:scripts}
% ══════════════════════════════════════════════════════════════════════════════

\subsection{Primary Script: \texttt{evaluation/run\_fair\_evaluation.py}}

This is the \textbf{single script} that orchestrates the full comparison.
Its 5-step pipeline:

\begin{enumerate}
    \item \textbf{Build stratified sample.}
          \texttt{tarbench\_bridge.build\_evaluation\_dataset()} loads
          TaRBench's \texttt{splits.csv} and all per-project
          \texttt{dataset.json} files, then draws a stratified sample.
    \item \textbf{Run GATeR.}
          For each sampled case, converts TaRBench metadata into GATeR's
          input format (\texttt{tarbench\_to\_gater\_input()}) and calls
          \texttt{GATREngine.repair\_test()}.
          Changed lines are extracted by diffing original vs.\ repaired code.
    \item \textbf{Load TARGET predictions.}
          Reads the pre-computed \texttt{test\_predictions.json} and
          \texttt{test\_verdicts.json} for the same sampled IDs.
    \item \textbf{Compute metrics.}
          Calls \texttt{metrics.compute\_all\_metrics()} for both systems.
    \item \textbf{Generate report.}
          Produces a Markdown comparison report and JSON results.
\end{enumerate}

\begin{lstlisting}[language=bash,caption={Running the evaluation}]
# Full evaluation (GATeR + TARGET on 500 cases, needs Ollama):
python -m evaluation.run_fair_evaluation

# TARGET-only (no Ollama needed, GATeR skipped):
python -m evaluation.run_fair_evaluation --skip-gater-run

# Custom sample size:
python -m evaluation.run_fair_evaluation --sample-size 200 --seed 42
\end{lstlisting}

\subsection{Supporting Modules}

\begin{table}[h]
\centering
\caption{Evaluation file inventory}
\begin{tabular}{@{}lp{8.5cm}@{}}
\toprule
\textbf{File} & \textbf{Purpose} \\
\midrule
\texttt{run\_fair\_evaluation.py}
    & Main 5-step evaluation driver; runs GATeR's full pipeline
      and loads TARGET's pre-computed results. \\
\texttt{metrics.py}
    & Shared metric functions: BLEU, CodeBLEU, EM, beam EM,
      plausible rate, compilation check, edit similarity,
      normalization, comparison table formatting. \\
\texttt{tarbench\_bridge.py}
    & TaRBench data loading, stratified sampling, TaRBench $\to$
      GATeR input conversion, ground-truth extraction, TARGET
      prediction loading. \\
\texttt{fair\_eval\_config.py}
    & All paths, sampling parameters (\texttt{SAMPLE\_SIZE=500},
      \texttt{RANDOM\_SEED=42}), metric configuration. \\
\bottomrule
\end{tabular}
\end{table}

% ══════════════════════════════════════════════════════════════════════════════
\section{Why This Comparison Is Fair}
\label{sec:fairness}
% ══════════════════════════════════════════════════════════════════════════════

\subsection{Same Task, Same Input Data}

Each TaRBench test case provides:
\begin{description}[style=nextline]
    \item[Test code] The broken Java test method (\texttt{bSource.code}).
    \item[Broken lines] The lines that must change (\texttt{hunk.sourceChanges}).
    \item[SUT changes] Source-under-test modifications that caused the
          breakage (\texttt{prioritized\_changes}).
    \item[Hunk type] Whether the repair is an addition, deletion, or
          modification.
    \item[Verdict] Whether the original break was a compilation error
          or a runtime test failure.
\end{description}

TARGET's encoder concatenates these fields into a tokenised sequence.
GATeR's bridge (\texttt{tarbench\_to\_gater\_input}) converts them into
the \texttt{broken\_test} dict and \texttt{error\_message} string that
\texttt{GATREngine} expects, then the pipeline retrieves \emph{additional}
context from the Knowledge Graph and vector store.

\subsection{Same Subset}

Both systems are evaluated on the exact same set of case IDs, drawn
via \texttt{stratified\_sample()} with:
\begin{itemize}
    \item Proportional representation by verdict type (compile error
          vs.\ test failure) and triviality flag.
    \item Per-project cap (\texttt{MAX\_CASES\_PER\_PROJECT=80}).
    \item Minimum project diversity ($\geq$10 projects).
    \item Fixed random seed (42) for full reproducibility.
\end{itemize}

The \texttt{sample\_ids\_<timestamp>.json} file is saved for audit.

\subsection{Same Normalization}

Both TARGET's tokenised output (e.g.,\ \texttt{method ( arg )}) and
GATeR's raw output (\texttt{method(arg)}) pass through the same
\texttt{normalize\_code()} function before any metric is computed:

\begin{lstlisting}[language=Python,caption={Code normalization}]
def normalize_code(code: str) -> str:
    s = " ".join(code.split())
    s = re.sub(r"\s*\(\s*", "(", s)
    s = re.sub(r"\s*\)\s*", ") ", s)
    s = re.sub(r"\s*\.\s*", ".", s)
    s = re.sub(r"\s*,\s*", ", ", s)
    s = re.sub(r"\s*;\s*", "; ", s)
    s = re.sub(r"\s*<\s*", "<", s)
    s = re.sub(r"\s*>\s*", "> ", s)
    return s.strip()
\end{lstlisting}

\subsection{Same Metrics}

The identical \texttt{compute\_all\_metrics()} function is called for
both systems. Beam EM and plausible rate are additionally reported for
TARGET as reference metrics (GATeR does not produce beam candidates or
execution verdicts in this pipeline).

% ══════════════════════════════════════════════════════════════════════════════
\section{Metrics Definitions}
\label{sec:metrics}
% ══════════════════════════════════════════════════════════════════════════════

\begin{table}[h]
\centering
\caption{Evaluation metrics by category}
\renewcommand{\arraystretch}{1.25}
\begin{tabular}{@{}p{3.4cm}p{6.8cm}c@{}}
\toprule
\textbf{Metric} & \textbf{Definition} & \textbf{Category} \\
\midrule
BLEU
    & Corpus-level 4-gram overlap between normalised predictions and
      ground truth, with smoothing (Method~1).
    & Fair \\
CodeBLEU
    & Code-aware metric: n-gram match, weighted keyword match,
      syntax match (AST sub-trees), and dataflow match, equally weighted.
      Uses TARGET's bundled implementation.
    & Fair \\
Compilation \%
    & Lightweight Java syntax check (balanced braces and parentheses).
      Not a real compiler; catches obvious structural failures.
    & Fair \\
Plausible Rate
    & Percentage of cases where at least one candidate compiles and
      passes the test (from TARGET's execution verdicts).
    & Fair \\
\midrule
Exact Match
    & Prediction == ground truth (after normalisation).
    & Biased\textsuperscript{\dag} \\
Beam EM ($k{=}40$)
    & Any of TARGET's 40 beam candidates matches.
    & Biased\textsuperscript{\dag} \\
\midrule
Edit Similarity
    & Average normalised Levenshtein ratio, 0--100.
    & Analysis \\
Line-Level Accuracy
    & Proportion of changed lines matched per case.
    & Analysis \\
Repair Attempt \%
    & Percentage of cases with a non-empty prediction.
    & Diagnostic \\
\bottomrule
\end{tabular}

\medskip\footnotesize
\textsuperscript{\dag} EM inherently favours fine-tuned models:
TARGET is trained to reproduce the exact ground-truth token sequence
and gets 40 beam attempts. GATeR generates via reasoning and may
produce \emph{functionally correct but syntactically different} repairs.
\end{table}

\subsection{Mathematical Formulations}

\textbf{BLEU} (corpus-level, smoothed):
\begin{equation}
    \text{BLEU} = \text{BP} \cdot \exp\!\left(\sum_{n=1}^{4} w_n \log p_n\right)
\end{equation}

\textbf{KGCompass} (GATeR's context scoring, not a repair metric but
central to context quality):
\begin{equation}
    S(f) = \beta^{\,l(f)} \;\Big[\;
        \alpha \cdot \cos(\mathbf{e}_i, \mathbf{e}_f)
        + (1{-}\alpha) \cdot \text{lev}(t_i, t_f)
    \;\Big]
    \qquad \alpha{=}0.3,\; \beta{=}0.6
\end{equation}

\textbf{Edit Similarity} per case:
\begin{equation}
    \text{EditSim}(t,p) = \frac{2 \cdot |LCS(t, p)|}{|t| + |p|}
\end{equation}

% ══════════════════════════════════════════════════════════════════════════════
\section{Paradigm Comparison}
% ══════════════════════════════════════════════════════════════════════════════

\begin{table}[h]
\centering
\caption{Paradigm differences between GATeR and TARGET}
\renewcommand{\arraystretch}{1.2}
\begin{tabular}{@{}p{3.6cm}p{5.2cm}p{5.2cm}@{}}
\toprule
\textbf{Dimension} & \textbf{GATeR} & \textbf{TARGET} \\
\midrule
Paradigm & RAG + zero-shot generation & Supervised fine-tuning \\
Training data & None & 36,639 repair pairs \\
Base model & DeepSeek Coder 6.7B & CodeT5+ 770M \\
Context retrieval & KG (Kuzu) + Vector store (LanceDB)
                    + KGCompass scoring
                  & SUT changes appended to input \\
Context processing & 6-step compression + 4-step RAG
                     aggregation
                   & Tokenised concatenation \\
Repair generation & LLM prompt ($T{=}0.1$) & Beam search ($k{=}40$) \\
Repository access & Yes (KG built from full repo) & No \\
Generalises to new projects & Yes---no retraining & Requires retraining \\
Latency per case & $\sim$10--60\,s & $\sim$0.5\,s (GPU) \\
\bottomrule
\end{tabular}
\end{table}

% ══════════════════════════════════════════════════════════════════════════════
\section{Step-by-Step Evaluation Workflow for FYP Defense}
\label{sec:workflow}
% ══════════════════════════════════════════════════════════════════════════════

\begin{enumerate}[label=\textbf{Step \arabic*:},leftmargin=3.5em]
    \item \textbf{Prepare Knowledge Graphs} \\
          For each TaRBench project in the sample, clone the repository
          at the correct commit and run GATeR's parser to build the Kuzu
          KG and LanceDB vector index.

    \item \textbf{Start Ollama} \\
          \texttt{ollama serve} \\
          \texttt{ollama list} (verify \texttt{deepseek-coder:6.7b}
          is present).

    \item \textbf{Run Evaluation} \\
          \texttt{python -m evaluation.run\_fair\_evaluation --sample-size 500}\\
          This runs GATeR's \emph{full pipeline} on each case and loads
          TARGET's pre-computed predictions for the same subset.

    \item \textbf{(Optional) TARGET-Only Quick Check} \\
          \texttt{python -m evaluation.run\_fair\_evaluation --skip-gater-run}\\
          Validates TARGET's metrics on the sample without needing Ollama.

    \item \textbf{Inspect Report} \\
          Outputs in \texttt{evaluation/results/}:
          \begin{itemize}
              \item \texttt{fair\_evaluation\_report\_<ts>.md} ---
                    Markdown comparison report.
              \item \texttt{fair\_evaluation\_results\_<ts>.json} ---
                    Raw metric values.
              \item \texttt{sample\_ids\_<ts>.json} ---
                    IDs used (for reproducibility).
          \end{itemize}

    \item \textbf{Validated Reference: TARGET on 500-Case Sample} \\
          TARGET's subset performance has been pre-validated:
          BLEU\,=\,62.37, EM\,=\,62.40\%, Plausible\,=\,81.20\%
          (vs.\ full-set EM\,=\,66.08\%, Plausible\,=\,80.04\%).
          This confirms the stratified sample is representative.
\end{enumerate}

% ══════════════════════════════════════════════════════════════════════════════
\section{Why Exact Match Alone Is Insufficient}
% ══════════════════════════════════════════════════════════════════════════════

Exact Match is the headline metric in the TARGET paper, but it is
structurally biased toward fine-tuned models:

\begin{itemize}
    \item \textbf{TARGET} is trained to reproduce the \emph{exact tokenised
          ground truth}. With 40 beams it gets 40 chances.
    \item \textbf{GATeR} generates code through \emph{KG-guided reasoning}.
          It may produce functionally correct repairs that differ
          syntactically.
\end{itemize}

\textbf{Example:}
\begin{lstlisting}[language=Java]
// Ground truth (tokenised):
assertEquals ( "test" , result . getName ( ) )
// GATeR output (raw):
Assert.assertEquals("test", result.getName())
\end{lstlisting}

These are semantically identical, but Exact Match ${}= 0$.
BLEU, CodeBLEU, and edit similarity correctly capture the near-perfect
overlap. The primary comparison should therefore use these
\emph{degree-of-correctness} metrics rather than binary EM.

% ══════════════════════════════════════════════════════════════════════════════
\section{Configuration \& Environment}
% ══════════════════════════════════════════════════════════════════════════════

\begin{table}[h]
\centering
\caption{Key paths and parameters}
\renewcommand{\arraystretch}{1.15}
\begin{tabular}{@{}lp{8cm}@{}}
\toprule
\textbf{Parameter} & \textbf{Default / Location} \\
\midrule
\texttt{TARBENCH\_ROOT} & \texttt{d:\textbackslash ...\textbackslash TaRBench\textbackslash TaRBench} \\
\texttt{TARGET\_RESULTS\_DIR} & \texttt{d:\textbackslash ...\textbackslash Best\_on\_TaRBench} \\
\texttt{TARGET\_FINETUNING\_DIR} & \texttt{d:\textbackslash ...\textbackslash fine-tuning}
    (needed for CodeBLEU keywords) \\
\texttt{OLLAMA\_BASE\_URL} & \texttt{http://localhost:11434} \\
\texttt{OLLAMA\_MODEL} & \texttt{deepseek-coder:6.7b} \\
\texttt{SAMPLE\_SIZE} & 500 \\
\texttt{RANDOM\_SEED} & 42 \\
\bottomrule
\end{tabular}
\end{table}

% ══════════════════════════════════════════════════════════════════════════════
\section{What GATeR's Pipeline Adds Beyond Raw LLM Prompting}
\label{sec:pipeline-value}
% ══════════════════════════════════════════════════════════════════════════════

An evaluation that bypasses GATeR's pipeline and simply feeds TaRBench
context into the LLM would be testing ``prompt engineering + LLM,'' not
GATeR. The following components are essential to GATeR's contribution
and \emph{must} be exercised during evaluation:

\begin{enumerate}
    \item \textbf{Knowledge Graph (Kuzu)} --- builds entity--relationship
          structure over the entire project (classes, methods, imports,
          CALLS, TESTS, MODIFIES edges).
    \item \textbf{KGCompass Relevance Scoring} --- scores candidate
          entities using graph-distance-weighted similarity
          (Eq.~$S(f)=\beta^{l(f)}[\alpha\cos + (1{-}\alpha)\text{lev}]$).
    \item \textbf{Context Compression (Steps 2.1--2.6)} --- hybrid
          scoring, threshold filtering, snippet compression, relationship
          pruning, pattern extraction, semantic examples.
    \item \textbf{RAG Aggregation (Steps 3.1--3.4)} --- entity clustering,
          API delta extraction, canonical usage identification, repair
          strategy formulation.
    \item \textbf{Structured LLM Prompt} --- the compressed, aggregated
          context is fed to DeepSeek Coder, not the raw TaRBench fields.
\end{enumerate}

\texttt{run\_fair\_evaluation.py} calls \texttt{GATREngine.repair\_test()}
which executes all five components. This is the only script that
evaluates the real GATeR system.

% ══════════════════════════════════════════════════════════════════════════════
\section{How TARGET's Metrics Are Computed from Pre-Computed JSON}
\label{sec:target-metrics}
% ══════════════════════════════════════════════════════════════════════════════

TARGET does not need a running model at evaluation time.
Its two JSON files contain everything needed to compute all metrics
on the sampled subset. The script works as follows:

\begin{enumerate}
    \item \textbf{\texttt{test\_predictions.json}} is loaded.
          It contains 7,103 entries, each with:
          \begin{itemize}
              \item \texttt{target} --- the ground-truth repair line(s).
              \item \texttt{preds} --- a list of 40 beam candidates.
          \end{itemize}
    \item \textbf{\texttt{test\_verdicts.json}} is loaded.
          It contains 284,120 execution results (7,103~$\times$~40 beams
          minus empties). Each records whether that candidate
          \emph{actually compiled and passed} the test.
    \item The bridge function \texttt{load\_target\_predictions(ids)}
          filters to the sampled subset and picks the
          \textbf{best prediction per case}: the beam that matches the
          ground truth exactly (if any), otherwise beam~0.
    \item The same \texttt{compute\_all\_metrics()} function that
          processes GATeR's outputs is called with TARGET's
          predictions and verdicts.
\end{enumerate}

This is how each of the 6 shared metrics gets its TARGET value:

\begin{table}[h]
\centering
\caption{How each metric is computed for TARGET from JSON files}
\renewcommand{\arraystretch}{1.3}
\begin{tabular}{@{}p{3.2cm}p{10cm}@{}}
\toprule
\textbf{Metric} & \textbf{How TARGET's value is obtained} \\
\midrule
BLEU
    & \texttt{best\_pred} (best beam candidate) for each sampled case is
      normalised and compared to the ground-truth repair lines using
      corpus-level BLEU with smoothing (Method~1). \\
CodeBLEU
    & Same \texttt{best\_pred} vs.\ ground truth, passed through
      TARGET's own CodeBLEU implementation (n-gram + keyword +
      syntax-tree + dataflow match). \\
Compilation~\%
    & \texttt{best\_pred} is checked for balanced braces and
      parentheses via \texttt{check\_java\_syntax()}. \\
Edit Similarity
    & Normalised Levenshtein ratio between the normalised
      \texttt{best\_pred} and normalised ground truth, averaged. \\
Line-Level Accuracy
    & Each changed line in the ground truth is checked for presence
      in the prediction (after normalisation); results averaged. \\
Repair Attempt~\%
    & Percentage of sampled cases where \texttt{best\_pred} is
      non-empty (TARGET nearly always produces output, so this is
      close to 100\%). \\
\bottomrule
\end{tabular}
\end{table}

\textbf{Key point}: No model inference is performed. TARGET's 6 shared
metrics are recomputed \emph{from its saved predictions} using exactly
the same code path as GATeR. Only the plausible rate and beam EM
require TARGET-specific data (execution verdicts and 40 beams).

To get TARGET's 6 metric scores \emph{without running GATeR}, use:
\begin{lstlisting}[language=bash]
python -m evaluation.run_fair_evaluation --skip-gater-run
\end{lstlisting}
This takes under 1 minute and needs no GPU or Ollama.

% ══════════════════════════════════════════════════════════════════════════════
\section{FAQ: Metric Coverage and Plausible Rate}
\label{sec:faq}
% ══════════════════════════════════════════════════════════════════════════════

\subsection{Which metrics are computed for both systems?}

All 6 primary metrics are computed for \textbf{both} GATeR and TARGET
through the \emph{exact same} function call
(\texttt{compute\_all\_metrics()}):

\begin{table}[h]
\centering
\caption{Metric coverage per system}
\renewcommand{\arraystretch}{1.2}
\begin{tabular}{@{}lccc@{}}
\toprule
\textbf{Metric} & \textbf{GATeR} & \textbf{TARGET} & \textbf{Same code path?} \\
\midrule
BLEU              & \checkmark & \checkmark & Yes \\
CodeBLEU          & \checkmark & \checkmark & Yes \\
Compilation~\%    & \checkmark & \checkmark & Yes \\
Edit Similarity   & \checkmark & \checkmark & Yes \\
Line-Level Acc.   & \checkmark & \checkmark & Yes \\
Repair Attempt~\% & \checkmark & \checkmark & Yes \\
Exact Match       & \checkmark & \checkmark & Yes (flagged biased) \\
\midrule
Plausible Rate    & N/A        & \checkmark & No --- see below \\
Beam EM ($k{=}40$) & N/A       & \checkmark & No --- TARGET-only \\
\bottomrule
\end{tabular}
\end{table}

\subsection{Why is Plausible Rate N/A for GATeR?}

\textbf{Plausible Rate} means: ``Did the repaired test actually compile
and pass when executed against the real project?''

TARGET has this metric because the TaRGET authors \emph{physically ran}
each of the 40 beam candidates in each project's build environment.
The results are stored in \texttt{test\_verdicts.json}
(284,120 execution verdicts).

GATeR does \textbf{not} have plausible rate because obtaining it would
require:
\begin{enumerate}
    \item Checking out each of the 59 projects at the exact breakage commit.
    \item Applying GATeR's generated repair to the broken test file.
    \item Running the project's full build (Maven/Gradle) and executing
          the specific test.
    \item Recording pass/fail.
\end{enumerate}

This is a \textbf{major infrastructure task} (different Java versions,
different build tool versions, different dependency trees per project)
that is not currently implemented in GATeR's evaluation pipeline.
The metric is therefore honestly reported as N/A rather than using
a misleading proxy.

\textbf{Note}: GATeR's \texttt{repair\_test()} returns a
\texttt{success} flag, but this only indicates whether the pipeline
produced output---it is \emph{not} an execution verdict.
The evaluation correctly passes \texttt{verdicts=None} for GATeR so
that plausible rate shows as N/A.

\subsection{Why is Beam EM N/A for GATeR?}

TARGET generates 40 beam candidates per case and reports Exact Match
if \emph{any} beam matches. GATeR generates a single repair (one LLM
call per case). Comparing a single attempt against 40 attempts would
not be meaningful, so beam EM is reported only for TARGET.

% ══════════════════════════════════════════════════════════════════════════════
\section{Summary}
% ══════════════════════════════════════════════════════════════════════════════

\begin{itemize}
    \item \textbf{One script} (\texttt{run\_fair\_evaluation.py}) evaluates
          both systems on the same stratified subset.
    \item \textbf{GATeR's full pipeline} (KG + KGCompass + compression +
          RAG + LLM) is executed for every case.
          TARGET's pre-computed predictions are loaded from JSON.
    \item \textbf{6 metrics} (BLEU, CodeBLEU, Compilation~\%, Edit
          Similarity, Line-Level Accuracy, Repair Attempt~\%) plus
          Exact Match are computed for \textbf{both} systems through
          the same code path.
    \item \textbf{Plausible Rate} is TARGET-only (requires execution
          verdicts GATeR does not produce).
    \item \textbf{Beam EM ($k{=}40$)} is TARGET-only (GATeR generates
          one repair, not 40).
    \item \textbf{Fairness} is ensured by: same task, same subset,
          same normalization, same metric implementations.
    \item \textbf{TARGET-only quick check}: run with
          \texttt{-{}-skip-gater-run} to get TARGET's 6 metrics in
          under 1 minute---no GPU, no Ollama.
    \item \textbf{GATeR's advantage}: zero-shot generalisation---no
          training data, no retraining, works on any Java project.
\end{itemize}

\end{document}
